\subsection{Logistic Regression Model}

One of the approaches taken by a participant is finetuning a DeBERTa (deberta-vs-base) model. 
They finetuned on 90\% of the given data and used the remaining 10\% for validation across 10 epochs. 
They achieved an accuracy of 0.62 and a (weighted) f1 score of 0.59 on the validation set. 
On the other hand, when testing the model on the provided test set, they achieved a (weighted) f1 score of 0.51 and an accuracy of 0.53. 
These results were expected because, after constructing a couple of models, and trying different variations for the task, when using DistilBERT word embeddings with a logistic regression model and using it to predict the genres, the same results were obtained for the test set. 

Another approach taken by another participant is finetuning a RoBERTa (Roberta-base) model. 
As can be seen below, the model scored a f1-score of 0.53 and an accuracy of 0.55. 
This surprisingly performed better than both the DeBERTa and logistic regression DistilBERT models. 
This may be due to the tokens being more specific to certain contexts or length of the movie plots. 
Furthermore, this may also be due to the splitting of the data, optimization of the model parameters, etc.
